\documentclass[../main.tex]{subfiles}
\ifSubfilesClassLoaded{\setcounter{chapter}{0}}{}
\begin{document}

\chapter{Pendahuluan, Identitas Modul, dan Pemetaan OBE}

\section{Identitas Mata Kuliah Praktikum}

Bab ini menjadi peta navigasi praktikum: bagaimana \textbf{Sub-CPMK P1--P10}, \textbf{CPMK-P1--P4}, dan \textbf{CPL prodi} dirangkai dalam 16 minggu. Mengetahui identitas MK menyelaraskan ekspektasi dosen, asisten, dan mahasiswa terhadap bukti kinerja di notebook.

Praktikum bersifat \textit{hands-on} satu SKS; fokus utamanya adalah kemampuan mengoperasikan rantai kerja data---bukan sekadar menghafal definisi. Artefak yang dinilai adalah notebook yang dapat direproduksi, dilengkapi sumber data, dan narasi keputusan analitik.

\begin{description}[style=nextline, widest=Dosen pengampu / koordinator praktikum:, itemsep=0.35em]
  \item[Nama mata kuliah:] Praktikum Data Science
  \item[Kode:] TIF-xxx-prak
  \item[Bobot:] 1 SKS (100\% praktikum; P = 1; setara $\sim$170 menit terstruktur per minggu)
  \item[Program studi:] Teknik Informatika (S1) --- Universitas Bale Bandung
  \item[Semester / tahun akademik:] \ldots
  \item[Dosen pengampu / koordinator praktikum:] \ldots
  \item[Prasyarat / ko-requisit:] MK \textbf{Data Science} (teori) bersamaan atau telah lulus (atau setara)
\end{description}

\noindent\textbf{Versi modul:} 2.0 \quad \textbf{Revisi terakhir:} 2026

\noindent\textbf{Artefak latihan:} folder \path{notebooks/minggu_01.ipynb} s.d.\ \path{notebooks/minggu_16.ipynb}.

\begin{table}[htbp]
\centering
\caption{Ringkasan artefak utama praktikum.}
\label{tab:artefak-modul}
\begin{tabular}{@{}p{3.2cm}p{9.5cm}@{}}
\toprule
\textbf{Artefak} & \textbf{Keterangan} \\
\midrule
Notebook mingguan & Bukti pengerjaan per Sub-CPMK; wajib lulus \textit{Restart Kernel and Run All}. \\
PDF modul & Panduan teori singkat, prosedur lab, dan contoh kode rujukan. \\
Silabus/RPS & Acuan CPL/CPMK dan alokasi waktu resmi MK. \\
\bottomrule
\end{tabular}
\end{table}

\section{Deskripsi dan Ruang Lingkup}

Modul ini mendukung \textbf{CPMK-P1--P4} dan \textbf{Sub-CPMK-P1--P10} untuk praktikum berbasis \textbf{OBE}. Lingkup: Python 3, \jupyter{} (atau Colab/VS Code), NumPy, Pandas, Matplotlib, Seaborn, \sklearn{}. Alur kerja: pemahaman masalah (CRISP-DM), data, prapemrosesan, EDA, pemodelan, evaluasi, komunikasi \textit{insight}. RPS resmi: \path{silabus_praktikum_data_science.tex}.

Setiap minggu memiliki \textbf{keluaran konkret} di notebook---dari etika dan kerangka CRISP-DM, manipulasi dan visualisasi data, hingga pemodelan dengan metrik yang dapat diinterpretasikan---sebagai jembatan antara teori Data Science dan praktik reproduktif di laboratorium.

\begin{figure}[htbp]
\centering
\begin{tikzpicture}[node distance=5mm and 4mm]
  \node[dsbox, fill=blue!12] (n1) {Business\\Und.};
  \node[dsbox, fill=blue!12, right=of n1] (n2) {Data\\Und.};
  \node[dsbox, fill=blue!12, right=of n2] (n3) {Data\\Prep.};
  \node[dsbox, fill=blue!12, below=of n3] (n4) {Modeling};
  \node[dsbox, fill=blue!12, left=of n4] (n5) {Evaluation};
  \node[dsbox, fill=blue!12, left=of n5] (n6) {Deploy};
  \draw[dsarrow] (n1) -- (n2);
  \draw[dsarrow] (n2) -- (n3);
  \draw[dsarrow] (n3) -- (n4);
  \draw[dsarrow] (n4) -- (n5);
  \draw[dsarrow] (n5) -- (n6);
\end{tikzpicture}
\caption{Enam fase CRISP-DM sebagai kerangka berpikir proyek data. Pada praktikum 1 SKS, notebook tidak wajib menampilkan semua fase penuh setiap minggu, tetapi narasi sebaiknya menunjukkan fase mana yang sedang dikerjakan.}
\end{figure}


\section{CPL Prodi yang Dibebankan}

\textbf{CPL} (\textit{Capaian Pembelajaran Lulusan}) menyatakan kemampuan tingkat prodi yang dibuktikan secara tidak langsung melalui CPMK dan Sub-CPMK praktikum. Dua CPL berikut menjadi jangkar penilaian proyek dan tugas akhir praktikum.

Pada praktikum, bukti CPL muncul dari kombinasi: (1) analisis masalah dan pemilihan metode yang masuk akal, (2) implementasi kode yang terbaca dan dapat diuji ulang, serta (3) komunikasi temuan kepada audiens non-teknis dalam bentuk ringkasan Markdown.

\begin{itemize}[leftmargin=*]
  \item \textbf{CPL01:} Menganalisis persoalan \textit{computing} kompleks dan menerapkan prinsip \textit{computing} serta disiplin relevan untuk mengidentifikasi solusi, dengan wawasan transdisiplin.
  \item \textbf{CPL03:} Merancang dan menganalisis algoritma untuk permasalahan organisasi, memilih dan menerapkannya pada bahasa pemrograman tertentu.
\end{itemize}

\begin{table}[htbp]
\centering
\caption{CPL yang dilibatkan dan contoh bukti di praktikum.}
\label{tab:cpl-bukti}
\begin{tabular}{@{}p{1.8cm}p{10.8cm}@{}}
\toprule
\textbf{CPL} & \textbf{Contoh bukti di notebook} \\
\midrule
CPL01 & Narasi pemahaman bisnis/data, keterbatasan studi, rekomendasi berbasis plot/metrik. \\
CPL03 & Pipeline \texttt{sklearn}, perbandingan algoritma, validasi silang, kode terstruktur. \\
\bottomrule
\end{tabular}
\end{table}

\section{CPMK Praktikum (kemampuan akhir MK)}

\textbf{CPMK-P1--P4} memecah tujuan MK menjadi empat blok kemampuan yang dapat dinilai di akhir semester. Setiap blok didukung oleh beberapa Sub-CPMK mingguan; penilaian formatif di tiap minggu sebaiknya merujuk indikator yang relevan dengan CPMK tersebut.

Urutan CPMK mengikuti alur kerja data science umum: infrastruktur dan proses (P1), data dan EDA (P2), pemodelan (P3), evaluasi dan komunikasi (P4). Mahasiswa diharapkan dapat menjelaskan ``mengapa'' di balik setiap langkah, bukan hanya menyalin kode.

\begin{itemize}[leftmargin=*]
  \item \textbf{CPMK-P1:} Lingkungan Python/\jupyter{}, K3, CRISP-DM ringkas.
  \item \textbf{CPMK-P2:} Muat data, prapemrosesan, EDA, visualisasi di notebook.
  \item \textbf{CPMK-P3:} Model ML (regresi, klasifikasi, ensemble ringkas, klastering) dengan \sklearn{}.
  \item \textbf{CPMK-P4:} Metrik, validasi silang, tuning sederhana, narasi \textit{insight}.
\end{itemize}

\begin{table}[htbp]
\centering
\caption{CPMK praktikum dan ranah utama bukti.}
\label{tab:cpmk-ranah}
\begin{tabular}{@{}llp{7.2cm}@{}}
\toprule
\textbf{CPMK} & \textbf{Ranah} & \textbf{Fokus penilaian} \\
\midrule
P1 & Lingkungan \& proses & Notebook rapi, K3, alur CRISP-DM \\
P2 & Data \& visual & Kualitas data, plot informatif \\
P3 & Pemodelan & API \texttt{fit/predict}, perbandingan model \\
P4 & Evaluasi \& narasi & Metrik, CV, ringkasan untuk pemangku kepentingan \\
\bottomrule
\end{tabular}
\end{table}

\section{Sub-CPMK Praktikum P1--P10 (indikator ringkas)}

\textbf{Sub-CPMK} memecah CPMK menjadi indikator operasional per minggu. Daftar berikut memudahkan Anda mencocokkan tugas lab dengan rubrik penilaian; detail prosedur dan contoh kode ada pada bab terkait.

Indikator P1--P10 saling mengisi: misalnya P3 (cleaning) mendukung P4 (EDA) dan P5 (split/pemodelan). Jika satu minggu tertinggal, dokumentasikan keterangan di notebook agar asisten dapat memberi umpan balik yang tepat.

\begin{itemize}[leftmargin=*]
  \item \textbf{P1} --- \jupyter{}, kernel, etika \& lisensi data
  \item \textbf{P2} --- Muat data; inspeksi Pandas
  \item \textbf{P3} --- Cleaning, encoding, scaling, fitur sederhana
  \item \textbf{P4} --- EDA; Seaborn
  \item \textbf{P5} --- Supervised vs unsupervised; \texttt{train\_test\_split}; API estimator
  \item \textbf{P6} --- Regresi linear; error prediksi
  \item \textbf{P7} --- Klasifikasi (logistik, K-NN, pohon)
  \item \textbf{P8} --- Random Forest; K-Means
  \item \textbf{P9} --- Metrik, confusion matrix, ROC-AUC opsional; CV; GridSearchCV
  \item \textbf{P10} --- Ringkasan eksekutif; pipeline end-to-end mini
\end{itemize}

\begin{table}[htbp]
\centering
\caption{Sub-CPMK P1--P10 dan kata kunci bukti di notebook.}
\label{tab:subcpmk-kunci}
\footnotesize
\begin{tabular}{@{}clp{6.8cm}@{}}
\toprule
\textbf{P} & \textbf{Topik} & \textbf{Kata kunci bukti} \\
\midrule
P1 & Jupyter \& etika & Kernel, lisensi data, CRISP-DM \\
P2 & Pandas I/O & \texttt{read\_*}, \texttt{info}, \texttt{describe} \\
P3 & Prapemrosesan & Missing, duplikat, encoding, scaler \\
P4 & EDA & Histogram, boxplot, korelasi \\
P5 & API sklearn & \texttt{train\_test\_split}, \texttt{fit} \\
P6 & Regresi & RMSE/MAE, baseline \\
P7 & Klasifikasi & Akurasi/F1, confusion matrix \\
P8 & Ensemble \& cluster & RF, K-Means, silhouette \\
P9 & Tuning & CV, GridSearch, ROC opsional \\
P10 & Komunikasi & Ringkasan eksekutif, pipeline \\
\bottomrule
\end{tabular}
\end{table}

\section{\texorpdfstring{Matriks Sub-CPMK $\rightarrow$ CPMK $\rightarrow$ CPL}{Matriks Sub-CPMK ke CPMK ke CPL}}

Matriks berikut menyatukan tiga tingkat capaian: dari indikator mingguan (Sub-CPMK) ke capaian MK (CPMK) hingga jejak ke CPL prodi. Gunakan tabel ini saat menyusun portofolio atau refleksi akhir semester.

\begin{table}[htbp]
\centering
\caption{Pemetaan Sub-CPMK ke CPMK dan CPL.}
\label{tab:matriks-obe}
\begin{tabular}{@{}lll@{}}
\toprule
\textbf{Sub-CPMK} & \textbf{CPMK} & \textbf{CPL} \\
\midrule
P1 & P1 & CPL01 \\
P2, P3, P4 & P2 & CPL01, CPL03 \\
P5, P6, P7, P8 & P3 & CPL03 \\
P9, P10 & P4 & CPL01, CPL03 \\
\bottomrule
\end{tabular}
\end{table}

\section{Sinkronisasi 16 Minggu dengan RPS (Bagian 10)}

Jadwal berikut selaras dengan rencana pembelajaran semester (RPS) bagian jadwal/topik lab. Angka menit mengacu pada sesi terstruktur; mahasiswa dapat meluangkan waktu tambahan di luar jam untuk eksplorasi dan perbaikan notebook.

Koordinasi dengan dosen pengampu teori disarankan agar konsep statistik atau ML yang relevan di kelas teori muncul pada minggu praktikum yang mendekatinya---tanpa mensyaratkan urutan yang kaku jika kurikulum bergeser.

{\footnotesize
\begingroup\setlength{\tabcolsep}{3pt}%
\begin{longtable}{|c|p{1.4cm}|>{\raggedright\arraybackslash}p{5.1cm}|c|}
\hline
\textbf{Mg} & \textbf{Sub} & \textbf{Topik lab} & \textbf{Wkt} \\
\hline
\endfirsthead
\hline
\textbf{Mg} & \textbf{Sub} & \textbf{Topik lab} & \textbf{Wkt} \\
\hline
\endhead
\hline
\endfoot
1 & P1 & Jupyter, CRISP-DM, etika data & 170' \\
\hline
2 & P2 & Pandas: CSV/Excel, head, info, describe & 170' \\
\hline
3 & P3 & Missing, duplikat, outlier & 170' \\
\hline
4 & P3 & Encoding, scaling, fitur & 170' \\
\hline
5 & P4 & EDA univariat/bivariat, Seaborn & 170' \\
\hline
6 & P5 & train\_test\_split, fit/predict & 170' \\
\hline
7 & P6 & Regresi linear, RMSE/MAE & 170' \\
\hline
8 & P7 & Klasifikasi: logistik, K-NN, pohon & 170' \\
\hline
9 & P8 & Random Forest vs baseline & 170' \\
\hline
10 & P8 & K-Means, elbow/silhouette & 170' \\
\hline
11 & P9 & Presisi, recall, F1; ROC-AUC opsional & 170' \\
\hline
12 & P9 & Cross-validation, GridSearchCV & 170' \\
\hline
13 & P10 & Narasi Markdown, ringkasan eksekutif & 170' \\
\hline
14 & P10 & Peer review, reproduktifitas & 170' \\
\hline
15 & P10 & Proyek mini end-to-end (draft) & 170' \\
\hline
16 & P10 & Presentasi \& penyerahan akhir & 170' \\
\hline
\end{longtable}%
\endgroup
}

\section{K3, Etika Data, dan Tata Tertib Lab}

K3 laboratorium melindungi pengguna perangkat dan menjaga keberlangsungan praktikum multi-angkatan; etika data memastikan penggunaan dataset sesuai lisensi dan norma privasi. Keduanya termasuk capaian profesionalisme, bukan sekadar administrasi.

Larangan makan/minum di workstation; laporkan kerusakan perangkat; patuhi lisensi dataset; jangan menyebarkan data pribadi tanpa izin; cantumkan atribusi tutorial dan sumber data di Markdown.

\begin{table}[htbp]
\centering
\caption{Ringkasan K3 dan etika data di lab.}
\label{tab:k3-etika}
\begin{tabular}{@{}p{3.5cm}p{9.2cm}@{}}
\toprule
\textbf{Aspek} & \textbf{Praktik yang diharapkan} \\
\midrule
K3 fisik & Tidak makan/minum di meja kerja; hati-hati kabel; laporkan perangkat panas. \\
Etika akademik & Kutip sumber kode/tutorial; tidak plagiasi notebook rekan. \\
Etika data & Cantumkan lisensi; hindari PII tanpa izin; hormati syarat API. \\
\bottomrule
\end{tabular}
\end{table}

\section{Petunjuk Penggunaan Modul}

Modul PDF dan notebook saling melengkapi: PDF memberi kerangka konsep dan urutan aktivitas, sedangkan notebook menjadi tempat eksperimen dan bukti capaian. Biasakan menulis Markdown antara sel kode agar pembaca lain (termasuk Anda di minggu depan) memahami maksud setiap langkah.

\begin{enumerate}[leftmargin=*]
  \item Baca \textbf{Informasi sesi praktikum} dan \textbf{Sub-CPMK} di tiap bab.
  \item Kerjakan \texttt{notebooks/minggu\_NN.ipynb} paralel dengan membaca bab PDF.
  \item Kumpulkan \texttt{.ipynb} yang lulus \textit{Restart Kernel and Run All}.
\end{enumerate}

\begin{figure}[htbp]
\centering
\begin{tikzpicture}[node distance=6mm]
  \node[dsboxwide, fill=purple!8] (m) {Bab PDF\\modul};
  \node[dsboxwide, fill=purple!8, right=of m] (n) {Notebook\\\texttt{.ipynb}};
  \node[dsboxwide, fill=purple!8, right=of n] (r) {Run All\\+ narasi};
  \draw[dsarrow] (m) -- (n);
  \draw[dsarrow] (n) -- (r);
\end{tikzpicture}
\caption{Alur belajar disarankan: baca tujuan sesi di PDF, lalu implementasikan dan dokumentasikan di notebook hingga seluruh sel dapat dijalankan ulang dari awal.}
\end{figure}


\begin{rangkuman}
Bab 0 memuat identitas MK, CPL/CPMK/Sub-CPMK P1--P10, jadwal 16 minggu selaras RPS, K3, etika data, dan alur notebook.

\textbf{Kata kunci:} OBE, CRISP-DM, \jupyter{}, TIF-xxx-prak
\end{rangkuman}

\end{document}
